\chapter{内积空间与谱定理}

在线性代数中，内积把“长度”“夹角”“垂直”这些几何概念引入到抽象向量空间之中。借助内积，我们不但能够构造正交基、建立最小二乘理论，而且还能把上一章关于特征值与特征向量的讨论推进到一个更深刻的层次：对于自伴算子与正规算子，可以找到一组正交归一基把它们完全对角化，这就是谱定理。对模形式而言，Petersson 内积正是把这些线性代数思想带到函数空间中的桥梁；Hecke 算子的自伴性与交换性，最终导向本征模形式分解。

本章默认向量空间是有限维的；当讨论复向量空间时，我们总采用“对第一变量线性、对第二变量共轭线性”的内积约定。

\section{内积空间的定义}

\subsection{实内积与复内积}

\begin{definition}
设 $V$ 是一个实向量空间。映射
\[
\inner{\cdot}{\cdot} : V \times V \longrightarrow \R
\]
称为 $V$ 上的一个\emph{实内积}，如果对任意 $u,v,w\in V$ 及任意 $a,b\in\R$，满足：
\begin{enumerate}[label=\textnormal{(\arabic*)}]
  \item 对称性：$\inner{u}{v}=\inner{v}{u}$；
  \item 双线性：$\inner{au+bv}{w}=a\inner{u}{w}+b\inner{v}{w}$，且同理对第二变量也线性；
  \item 正定性：$\inner{v}{v}\ge 0$，且 $\inner{v}{v}=0$ 当且仅当 $v=0$。
\end{enumerate}
称带有内积的实向量空间为\emph{实内积空间}。
\end{definition}

\begin{definition}
设 $V$ 是一个复向量空间。映射
\[
\inner{\cdot}{\cdot}:V\times V\longrightarrow \C
\]
称为 $V$ 上的一个\emph{Hermite 内积}或\emph{复内积}，如果对任意 $u,v,w\in V$ 及任意 $a,b\in\C$，满足：
\begin{enumerate}[label=\textnormal{(\arabic*)}]
  \item 共轭对称性：$\inner{u}{v}=\overline{\inner{v}{u}}$；
  \item 第一变量线性：$\inner{au+bv}{w}=a\inner{u}{w}+b\inner{v}{w}$；
  \item 第二变量共轭线性：$\inner{u}{av+bw}=\overline{a}\inner{u}{v}+\overline{b}\inner{u}{w}$；
  \item 正定性：$\inner{v}{v}\ge 0$，且 $\inner{v}{v}=0$ 当且仅当 $v=0$。
\end{enumerate}
称带有这种内积的复向量空间为\emph{复内积空间}。
\end{definition}

\begin{remark}
许多分析书把复内积取为“对第二变量线性”。两种约定完全等价，但一旦选定，就必须在全书中保持一致。本书采用对第一变量线性的约定，因为这样与矩阵左乘、算子作用的记号更为协调。
\end{remark}

\begin{example}
在 $\R^n$ 上定义
\[
\inner{x}{y}=x_1y_1+\cdots+x_ny_n.
\]
这就是标准实内积，也称点积。
\end{example}

\begin{example}
在 $\C^n$ 上定义
\[
\inner{x}{y}=x_1\overline{y_1}+\cdots+x_n\overline{y_n}.
\]
这就是标准 Hermite 内积。特别地，
\[
\inner{x}{x}=|x_1|^2+\cdots+|x_n|^2\ge 0,
\]
且等号仅当 $x=0$ 时成立。
\end{example}

\begin{example}
在 $\R^n$ 上，设 $A\in M_n(\R)$ 是对称正定矩阵，定义
\[
\inner{x}{y}_A=x^tAy.
\]
则它给出一个内积。对称性来自 $A^t=A$，双线性显然，而正定性正是 $x^tAx>0$（对一切非零 $x$）的含义。
\end{example}

\begin{example}
设 $P_{m}(\R)$ 表示次数不超过 $m$ 的实系数多项式空间。定义
\[
\inner{f}{g}=\int_0^1 f(x)g(x)\,dx.
\]
它是 $P_m(\R)$ 上的内积。这里正定性的关键在于：若 $\int_0^1 f(x)^2\,dx=0$，则连续函数 $f^2$ 在区间上恒为零，从而 $f=0$。
\end{example}

\begin{proposition}
设 $V$ 是内积空间，则对任意 $u,v,w\in V$ 有：
\begin{enumerate}[label=\textnormal{(\arabic*)}]
  \item $\inner{v}{0}=\inner{0}{v}=0$；
  \item $\inner{u}{v+w}=\inner{u}{v}+\inner{u}{w}$；
  \item $\inner{u}{v}=0$ 对一切 $u\in V$ 成立，当且仅当 $v=0$；
  \item $\inner{v}{v}\in\R$，且 $\inner{v}{v}\ge 0$。
\end{enumerate}
\end{proposition}

\begin{proof}
(1) 由线性可得
\[
\inner{v}{0}=\inner{v}{0+0}=\inner{v}{0}+\inner{v}{0},
\]
故 $\inner{v}{0}=0$。另一式同理。

(2) 在实情形这就是双线性；在复情形由第二变量的共轭线性可得
\[
\inner{u}{v+w}=\overline{1}\inner{u}{v}+\overline{1}\inner{u}{w}=\inner{u}{v}+\inner{u}{w}.
\]

(3) 若 $\inner{u}{v}=0$ 对一切 $u$ 成立，特别取 $u=v$，则 $\inner{v}{v}=0$，由正定性得 $v=0$。反过来若 $v=0$，由 (1) 立得。

(4) 在实情形显然。在复情形中，由共轭对称性，
\[
\inner{v}{v}=\overline{\inner{v}{v}},
\]
故 $\inner{v}{v}\in\R$；再由正定性知其非负。
\end{proof}

\subsection{范数与距离}

\begin{definition}
设 $V$ 是内积空间。对任意 $v\in V$，定义
\[
\|v\|=\sqrt{\inner{v}{v}},
\]
称为 $v$ 的\emph{范数}或\emph{长度}。进一步定义
\[
 d(u,v)=\|u-v\|,
\]
称为 $u,v$ 之间的\emph{距离}。

若 $\inner{u}{v}=0$，则称 $u$ 与 $v$ \emph{正交}，记作 $u\perp v$。
\end{definition}

\begin{proposition}[勾股恒等式]
若 $u\perp v$，则
\[
\|u+v\|^2=\|u\|^2+\|v\|^2.
\]
\end{proposition}

\begin{proof}
由内积展开可得
\[
\|u+v\|^2=\inner{u+v}{u+v}=\inner{u}{u}+\inner{u}{v}+\inner{v}{u}+\inner{v}{v}.
\]
因为 $u\perp v$，中间两项均为零，于是结论成立。
\end{proof}

\begin{theorem}[Cauchy--Schwarz 不等式]
设 $V$ 是内积空间，则对任意 $u,v\in V$，有
\[
|\inner{u}{v}|\le \|u\|\,\|v\|.
\]
且等号成立当且仅当 $u,v$ 线性相关。
\end{theorem}

\begin{proof}
若 $v=0$，结论显然。以下设 $v\ne 0$。对任意标量 $t$，由正定性有
\[
0\le \|u-tv\|^2=\inner{u-tv}{u-tv}.
\]
在复情形下展开得
\[
\|u-tv\|^2=\|u\|^2-t\inner{v}{u}-\overline{t}\inner{u}{v}+|t|^2\|v\|^2.
\]
取
\[
t=\frac{\inner{u}{v}}{\|v\|^2},
\]
便有
\[
0\le \|u\|^2-\frac{|\inner{u}{v}|^2}{\|v\|^2}.
\]
两边乘以 $\|v\|^2$，得
\[
|\inner{u}{v}|^2\le \|u\|^2\|v\|^2,
\]
于是
\[
|\inner{u}{v}|\le \|u\|\,\|v\|.
\]

现在讨论等号情形。若 $u$ 与 $v$ 线性相关，则设 $u=cv$，于是
\[
|\inner{u}{v}|=|c|\,\inner{v}{v}=\|u\|\,\|v\|,
\]
故等号成立。

反过来若等号成立，则上面的不等式成为等式，从而
\[
\|u-tv\|^2=0
\]
对 $t=\inner{u}{v}/\|v\|^2$ 成立。于是 $u=tv$，即 $u,v$ 线性相关。
\end{proof}

\begin{corollary}[三角不等式]
对任意 $u,v\in V$，有
\[
\|u+v\|\le \|u\|+\|v\|.
\]
\end{corollary}

\begin{proof}
由展开公式与 Cauchy--Schwarz 不等式，
\[
\|u+v\|^2=\|u\|^2+\inner{u}{v}+\inner{v}{u}+\|v\|^2
\le \|u\|^2+2|\inner{u}{v}|+\|v\|^2
\le (\|u\|+\|v\|)^2.
\]
两边取平方根即得结论。
\end{proof}

\begin{remark}
Cauchy--Schwarz 不等式是整个内积理论的基础。它立刻导出夹角公式
\[
\cos\theta=\frac{\inner{u}{v}}{\|u\|\,\|v\|}
\quad (u,v\ne 0),
\]
并保证右端绝对值不超过 $1$。
\end{remark}

\begin{example}
在 $\R^3$ 中，取 $u=(1,1,0)$，$v=(1,0,1)$。则
\[
\inner{u}{v}=1,\qquad \|u\|=\|v\|=\sqrt{2}.
\]
所以二者夹角满足
\[
\cos\theta=\frac{1}{2},
\]
故 $\theta=\pi/3$。
\end{example}

\section{正交性与 Gram--Schmidt 过程}

\subsection{正交集与正交归一集}

\begin{definition}
设 $V$ 是内积空间，$S=\{v_1,\dots,v_m\}\subseteq V$。
\begin{enumerate}[label=\textnormal{(\arabic*)}]
  \item 若当 $i\ne j$ 时总有 $\inner{v_i}{v_j}=0$，则称 $S$ 为\emph{正交集}；
  \item 若 $S$ 是正交集且每个 $v_i$ 都满足 $\|v_i\|=1$，则称 $S$ 为\emph{正交归一集}；
  \item 若 $S$ 张成整个空间 $V$，则称 $S$ 为 $V$ 的\emph{正交基}或\emph{正交归一基}。
\end{enumerate}
\end{definition}

\begin{proposition}
内积空间中任意一组两两正交的非零向量都是线性无关的。
\end{proposition}

\begin{proof}
设 $v_1,\dots,v_m$ 两两正交且均非零，并有
\[
a_1v_1+\cdots+a_mv_m=0.
\]
固定 $j$，将上式与 $v_j$ 取内积，得
\[
0=\inner{a_1v_1+\cdots+a_mv_m}{v_j}=a_j\inner{v_j}{v_j},
\]
因为当 $i\ne j$ 时 $\inner{v_i}{v_j}=0$。又因 $v_j\ne 0$，故 $\inner{v_j}{v_j}>0$，于是 $a_j=0$。对所有 $j$ 都成立，所以该组向量线性无关。
\end{proof}

\begin{proposition}
若 $e_1,\dots,e_m$ 是正交归一集，且 $v\in \\mathrm{span}\{e_1,\dots,e_m\}$，则
\[
v=\sum_{j=1}^m \inner{v}{e_j}e_j.
\]
\end{proposition}

\begin{proof}
设 $v=\sum_{j=1}^m a_je_j$。对任意固定的 $i$，与 $e_i$ 取内积，得
\[
\inner{v}{e_i}=\sum_{j=1}^m a_j\inner{e_j}{e_i}=a_i.
\]
故 $a_i=\inner{v}{e_i}$，从而得到所需展开式。
\end{proof}

\subsection{Gram--Schmidt 正交化}

\begin{theorem}[Gram--Schmidt 过程]
设 $v_1,\dots,v_m$ 是内积空间 $V$ 中的一组线性无关向量。定义
\[
u_1=v_1,
\]
并对 $k\ge 2$ 定义
\[
u_k=v_k-\sum_{j=1}^{k-1}\frac{\inner{v_k}{u_j}}{\inner{u_j}{u_j}}u_j.
\]
则有：
\begin{enumerate}[label=\textnormal{(\arabic*)}]
  \item 每个 $u_k$ 都非零；
  \item $u_1,\dots,u_m$ 两两正交；
  \item 对每个 $k$，有
  \[
  \\mathrm{span}\{u_1,\dots,u_k\}=\\mathrm{span}\{v_1,\dots,v_k\}.
  \]
\end{enumerate}
因此令
\[
e_k=\frac{u_k}{\|u_k\|}\qquad (1\le k\le m),
\]
则 $e_1,\dots,e_m$ 是与 $v_1,\dots,v_m$ 生成同一子空间的一组正交归一向量。
\end{theorem}

\begin{proof}
我们用归纳法证明。首先 $u_1=v_1\ne 0$。

对一般的 $k$，先证明正交性。若 $i<k$，则
\[
\inner{u_k}{u_i}
=\inner{v_k}{u_i}-\sum_{j=1}^{k-1}\frac{\inner{v_k}{u_j}}{\inner{u_j}{u_j}}\inner{u_j}{u_i}.
\]
由于归纳假设 $u_1,\dots,u_{k-1}$ 两两正交，上式中只有 $j=i$ 的一项可能不为零，因此
\[
\inner{u_k}{u_i}=\inner{v_k}{u_i}-\frac{\inner{v_k}{u_i}}{\inner{u_i}{u_i}}\inner{u_i}{u_i}=0.
\]
所以 $u_k$ 与前面的所有 $u_i$ 都正交。

其次证明 $u_k\ne 0$。若 $u_k=0$，则
\[
v_k=\sum_{j=1}^{k-1}\frac{\inner{v_k}{u_j}}{\inner{u_j}{u_j}}u_j\in \\mathrm{span}\{u_1,\dots,u_{k-1}\}.
\]
根据归纳假设，$\\mathrm{span}\{u_1,\dots,u_{k-1}\}=\\mathrm{span}\{v_1,\dots,v_{k-1}\}$，于是 $v_k$ 落在前 $k-1$ 个向量张成的空间中，与 $v_1,\dots,v_m$ 线性无关矛盾。所以 $u_k\ne 0$。

最后证明张成空间相同。由定义立刻有 $u_k\in \\mathrm{span}\{v_1,\dots,v_k\}$，故
\[
\\mathrm{span}\{u_1,\dots,u_k\}\subseteq \\mathrm{span}\{v_1,\dots,v_k\}.
\]
反过来，由
\[
v_k=u_k+\sum_{j=1}^{k-1}\frac{\inner{v_k}{u_j}}{\inner{u_j}{u_j}}u_j
\]
可知 $v_k\in \\mathrm{span}\{u_1,\dots,u_k\}$。结合归纳假设，得到反向包含，于是两者相等。

将每个非零向量 $u_k$ 归一化即可得到正交归一组 $e_1,\dots,e_m$，结论成立。
\end{proof}

\begin{corollary}
每个有限维内积空间都存在正交归一基。
\end{corollary}

\begin{proof}
任取一组基 $v_1,\dots,v_n$，对其施行 Gram--Schmidt 过程即可得到一组正交归一基。
\end{proof}

\begin{example}
在 $\R^3$ 中取
\[
v_1=\mat{1\\1\\0},\qquad v_2=\mat{1\\0\\1},\qquad v_3=\mat{0\\1\\1}.
\]
先取
\[
u_1=v_1=\mat{1\\1\\0},\qquad e_1=\frac{1}{\sqrt2}\mat{1\\1\\0}.
\]
再计算
\[
u_2=v_2-\frac{\inner{v_2}{u_1}}{\inner{u_1}{u_1}}u_1
=\mat{1\\0\\1}-\frac{1}{2}\mat{1\\1\\0}
=\mat{\tfrac12\\-\tfrac12\\1}.
\]
于是
\[
e_2=\frac{u_2}{\|u_2\|}=\frac{1}{\sqrt6}\mat{1\\-1\\2}.
\]
最后
\[
\inner{v_3}{e_1}=\frac{1}{\sqrt2},\qquad \inner{v_3}{e_2}=\frac{1}{\sqrt6},
\]
故
\[
u_3=v_3-\frac{1}{\sqrt2}e_1-\frac{1}{\sqrt6}e_2
=\mat{-\tfrac23\\\tfrac23\\\tfrac23}.
\]
归一化得
\[
e_3=\frac{1}{\sqrt3}\mat{-1\\1\\1}.
\]
因此 $e_1,e_2,e_3$ 构成 $\R^3$ 的一组正交归一基。
\end{example}

\subsection{QR 分解}

\begin{theorem}[QR 分解]
设 $A\in M_{m\times n}(\R)$ 的列向量为 $a_1,\dots,a_n$，并假设它们线性无关。则存在唯一的分解
\[
A=QR,
\]
其中
\begin{enumerate}[label=\textnormal{(\arabic*)}]
  \item $Q\in M_{m\times n}(\R)$ 的列向量 $q_1,\dots,q_n$ 正交归一；
  \item $R\in M_n(\R)$ 是上三角矩阵；
  \item $R$ 的对角元全为正数。
\end{enumerate}
\end{theorem}

\begin{proof}
对列向量 $a_1,\dots,a_n$ 应用 Gram--Schmidt 过程，得到正交归一向量 $q_1,\dots,q_n$，并且对每个 $j$ 有
\[
a_j=\sum_{i=1}^j r_{ij}q_i,
\]
其中
\[
r_{ij}=\inner{a_j}{q_i}\quad (i<j),\qquad r_{jj}=\|u_j\|>0.
\]
令 $Q=(q_1\ \cdots\ q_n)$，$R=(r_{ij})$，则由上式可知 $A=QR$，而 $R$ 显然是上三角矩阵。

下面证唯一性。若还有另一分解 $A=Q'R'$，其中 $Q'$ 也有正交归一列，$R'$ 也是对角元为正的上三角矩阵，则
\[
Q^TQ=I_n,\qquad (Q')^TQ'=I_n.
\]
由 $QR=Q'R'$ 得
\[
R=(Q^TQ')R'.
\]
记 $U=Q^TQ'$，则 $U$ 是正交矩阵。于是
\[
R=UR'.
\]
左边与右边都是上三角矩阵，所以 $U=RR'^{-1}$ 也是上三角矩阵。一个既正交又上三角的实矩阵必为对角矩阵，且对角元只能是 $\pm1$。又因为 $R,R'$ 的对角元都为正，故 $U$ 的对角元全为 $1$，从而 $U=I_n$。于是 $R=R'$，$Q=Q'$，唯一性得证。
\end{proof}

\begin{example}
取
\[
A=\mat{1&1\\1&0\\0&1}.
\]
其列向量为 $a_1=(1,1,0)^t$，$a_2=(1,0,1)^t$。由上面的计算可得
\[
q_1=\frac{1}{\sqrt2}\mat{1\\1\\0},\qquad q_2=\frac{1}{\sqrt6}\mat{1\\-1\\2}.
\]
于是
\[
Q=\mat{\frac1{\sqrt2}&\frac1{\sqrt6}\\[3pt]
\frac1{\sqrt2}&-\frac1{\sqrt6}\\[3pt]
0&\frac2{\sqrt6}},
\qquad
R=\mat{\sqrt2&\frac1{\sqrt2}\\0&\frac{\sqrt6}{2}}.
\]
直接相乘可验证 $A=QR$。
\end{example}

\section{正交补与正交投影}

\subsection{正交补}

\begin{definition}
设 $V$ 是内积空间，$U\subseteq V$ 是一个子空间。定义
\[
U^{\perp}=\{v\in V:\inner{v}{u}=0\text{ 对一切 }u\in U\text{ 都成立}\},
\]
称为 $U$ 的\emph{正交补}。
\end{definition}

\begin{proposition}
$U^{\perp}$ 是 $V$ 的子空间。
\end{proposition}

\begin{proof}
零向量显然属于 $U^{\perp}$。若 $v,w\in U^{\perp}$，$a,b$ 为标量，则对任意 $u\in U$，
\[
\inner{av+bw}{u}=a\inner{v}{u}+b\inner{w}{u}=0.
\]
故 $av+bw\in U^{\perp}$。因此 $U^{\perp}$ 是子空间。
\end{proof}

\begin{theorem}
设 $V$ 是有限维内积空间，$U\subseteq V$ 是子空间，则
\[
V=U\oplus U^{\perp}.
\]
特别地，每个 $v\in V$ 都可以唯一表示为
\[
v=u+w\qquad (u\in U,\ w\in U^{\perp}).
\]
\end{theorem}

\begin{proof}
先在 $U$ 中取一组正交归一基 $e_1,\dots,e_r$，再把它扩充为 $V$ 的一组正交归一基
\[
e_1,\dots,e_r,e_{r+1},\dots,e_n.
\]
令
\[
W=\\mathrm{span}\{e_{r+1},\dots,e_n\}.
\]
显然 $V=U\oplus W$。又对任意 $w\in W$ 与任意 $u\in U$，将 $u$ 写成 $e_1,\dots,e_r$ 的线性组合，则每一项都与 $w$ 正交，从而 $\inner{w}{u}=0$。故 $W\subseteq U^{\perp}$。

反过来，若 $x\in U^{\perp}$，则把 $x$ 按该正交归一基展开：
\[
x=\sum_{i=1}^n \inner{x}{e_i}e_i.
\]
对 $1\le i\le r$，因 $e_i\in U$ 且 $x\in U^{\perp}$，有 $\inner{x}{e_i}=0$。于是
\[
x=\sum_{i=r+1}^n \inner{x}{e_i}e_i\in W.
\]
故 $U^{\perp}\subseteq W$。于是 $W=U^{\perp}$，从而 $V=U\oplus U^{\perp}$。

唯一性来自直和分解：若 $u_1+w_1=u_2+w_2$，则 $(u_1-u_2)=(w_2-w_1)$ 同时属于 $U$ 与 $U^{\perp}$。于是
\[
\inner{u_1-u_2}{u_1-u_2}=0,
\]
故 $u_1=u_2$，再得 $w_1=w_2$。
\end{proof}

\begin{corollary}
在有限维内积空间中，对任意子空间 $U$，有
\[
(U^{\perp})^{\perp}=U.
\]
\end{corollary}

\begin{proof}
显然 $U\subseteq (U^{\perp})^{\perp}$。另一方面，由上一定理，
\[
V=U\oplus U^{\perp}
\]
且
\[
V=U^{\perp}\oplus (U^{\perp})^{\perp}.
\]
比较维数可得
\[
\dim (U^{\perp})^{\perp}=\dim V-\dim U^{\perp}=\dim U.
\]
由于已有包含关系，遂得两者相等。
\end{proof}

\subsection{正交投影}

\begin{definition}
设 $U\subseteq V$ 是子空间。对任意 $v\in V$，由分解定理可唯一写成
\[
v=u+w\qquad (u\in U,\ w\in U^{\perp}).
\]
称这个 $u$ 为 $v$ 在 $U$ 上的\emph{正交投影}，记作 $P_U(v)$。
\end{definition}

\begin{theorem}[投影公式]
设 $e_1,\dots,e_r$ 是子空间 $U$ 的一组正交归一基，则对任意 $v\in V$，
\[
P_U(v)=\sum_{i=1}^r \inner{v}{e_i}e_i.
\]
\end{theorem}

\begin{proof}
设
\[
u=\sum_{i=1}^r \inner{v}{e_i}e_i\in U.
\]
只需证明 $v-u\in U^{\perp}$。对任意 $j$，有
\[
\inner{v-u}{e_j}=\inner{v}{e_j}-\sum_{i=1}^r \inner{v}{e_i}\inner{e_i}{e_j}=\inner{v}{e_j}-\inner{v}{e_j}=0.
\]
故 $v-u$ 与 $U$ 的基向量都正交，于是 $v-u\in U^{\perp}$。由正交分解的唯一性知 $u=P_U(v)$。
\end{proof}

\begin{proposition}
正交投影 $P_U:V\to U$ 是线性映射，并满足
\[
P_U^2=P_U,\qquad \Ker(P_U)=U^{\perp}.
\]
\end{proposition}

\begin{proof}
由投影公式可见 $P_U$ 是线性的。又若 $v\in V$，则 $P_U(v)\in U$，所以再次投影时保持不变，即 $P_U(P_U(v))=P_U(v)$，从而 $P_U^2=P_U$。

若 $P_U(v)=0$，则在分解 $v=P_U(v)+(v-P_U(v))$ 中，第一项为零，故 $v\in U^{\perp}$。反之若 $v\in U^{\perp}$，则它在 $U$ 上的投影显然为零，所以 $\Ker(P_U)=U^{\perp}$。
\end{proof}

\begin{theorem}[最佳逼近定理]
设 $U\subseteq V$ 是子空间，$v\in V$。则 $u_0\in U$ 是 $v$ 在 $U$ 中的最佳逼近，即
\[
\|v-u_0\|\le \|v-u\|\qquad (\forall u\in U),
\]
当且仅当 $u_0=P_U(v)$。
\end{theorem}

\begin{proof}
先设 $u_0=P_U(v)$。则 $v-u_0\in U^{\perp}$。对任意 $u\in U$，有 $u_0-u\in U$，从而
\[
v-u=(v-u_0)+(u_0-u)
\]
是一个正交分解。于是由勾股恒等式，
\[
\|v-u\|^2=\|v-u_0\|^2+\|u_0-u\|^2\ge \|v-u_0\|^2.
\]
故 $u_0$ 的确是最佳逼近。

反过来，若 $u_0\in U$ 是最佳逼近。对任意 $u\in U$ 与任意实数 $t$，令
\[
\phi(t)=\|v-(u_0+tu)\|^2.
\]
由最佳逼近性质，$t=0$ 是 $\phi$ 的极小点。展开得
\[
\phi(t)=\|v-u_0\|^2-t\inner{u}{v-u_0}-t\inner{v-u_0}{u}+t^2\|u\|^2.
\]
在实情形下，对 $t$ 求导并令 $t=0$，得
\[
\inner{u}{v-u_0}+\inner{v-u_0}{u}=2\inner{u}{v-u_0}=0.
\]
在复情形中分别取 $u$ 与 $iu$，可推出 $\inner{u}{v-u_0}=0$。因此对一切 $u\in U$ 都有 $\inner{u}{v-u_0}=0$，即 $v-u_0\in U^{\perp}$。于是 $u_0=P_U(v)$。
\end{proof}

\begin{example}
在 $\R^3$ 中，设 $U=\\mathrm{span}\{(1,1,0)^t\}$，取 $v=(2,1,3)^t$。因为
\[
e=\frac{1}{\sqrt2}\mat{1\\1\\0}
\]
是 $U$ 的正交归一基，所以
\[
P_U(v)=\inner{v}{e}e=\frac{3}{\sqrt2}\cdot \frac{1}{\sqrt2}\mat{1\\1\\0}=\mat{\tfrac32\\\tfrac32\\0}.
\]
这就是 $v$ 在该直线上的最佳逼近。
\end{example}

\subsection{最小二乘法}

\begin{theorem}[最小二乘与正规方程]
设 $A\in M_{m\times n}(\C)$，$b\in\C^m$。则向量 $x_0\in\C^n$ 使得
\[
\|Ax_0-b\|\le \|Ax-b\|\qquad (\forall x\in\C^n)
\]
当且仅当它满足\emph{正规方程}
\[
A^*Ax_0=A^*b.
\]
若 $A$ 的列向量线性无关，则解唯一，并且
\[
x_0=(A^*A)^{-1}A^*b.
\]
\end{theorem}

\begin{proof}
设 $U=\im(A)\subseteq \C^m$。问题是在子空间 $U$ 中寻找离 $b$ 最近的向量，而所有形如 $Ax$ 的向量正好构成 $U$。因此 $x_0$ 是最小二乘解，当且仅当 $Ax_0=P_U(b)$。由最佳逼近定理，这又等价于
\[
b-Ax_0\in U^{\perp}.
\]

现在说明 $U^{\perp}=\Ker(A^*)$。若 $y\in U^{\perp}$，则对任意 $x\in\C^n$ 有
\[
0=\inner{Ax}{y}=\inner{x}{A^*y}.
\]
由内积的非退化性得 $A^*y=0$。反过来若 $A^*y=0$，则对任意 $x$，
\[
\inner{Ax}{y}=\inner{x}{A^*y}=0,
\]
所以 $y\in U^{\perp}$。故两者相等。

于是最小二乘条件等价于
\[
A^*(b-Ax_0)=0,
\]
即
\[
A^*Ax_0=A^*b.
\]

若 $A$ 的列向量线性无关，则 $A^*A$ 可逆。事实上对任意非零 $x$，
\[
\inner{A^*Ax}{x}=\inner{Ax}{Ax}=\|Ax\|^2>0,
\]
所以 $A^*A$ 正定，从而可逆。于是正规方程有唯一解，并给出所述公式。
\end{proof}

\begin{example}
考虑用直线 $y=\alpha+\beta t$ 拟合三点 $(0,1)$、$(1,2)$、$(2,2)$。令
\[
A=\mat{1&0\\1&1\\1&2},\qquad x=\mat{\alpha\\\beta},\qquad b=\mat{1\\2\\2}.
\]
则最小二乘解满足
\[
A^TAx=A^Tb.
\]
计算得
\[
A^TA=\mat{3&3\\3&5},\qquad A^Tb=\mat{5\\6}.
\]
解得
\[
x=\mat{\tfrac76\\\tfrac12}.
\]
因此最佳拟合直线为
\[
y=\frac76+\frac12 t.
\]
\end{example}

\section{自伴算子与谱定理}

\subsection{伴随算子}

\begin{theorem}
设 $V$ 是有限维内积空间，$T:V\to V$ 是线性算子。则存在唯一的线性算子 $T^*:V\to V$，使得对任意 $u,v\in V$ 都有
\[
\inner{Tu}{v}=\inner{u}{T^*v}.
\]
算子 $T^*$ 称为 $T$ 的\emph{伴随算子}。
\end{theorem}

\begin{proof}
先证存在性。取 $V$ 的一组正交归一基 $e_1,\dots,e_n$。固定 $v\in V$，定义线性泛函
\[
\varphi_v(u)=\inner{Tu}{v}.
\]
设
\[
w_v=\sum_{i=1}^n \overline{\varphi_v(e_i)}e_i.
\]
若 $u=\sum_i x_ie_i$，则
\[
\inner{u}{w_v}=\sum_i x_i\varphi_v(e_i)=\varphi_v(u)=\inner{Tu}{v}.
\]
因此存在向量 $w_v$ 使得 $\inner{Tu}{v}=\inner{u}{w_v}$ 对所有 $u$ 成立。令 $T^*v=w_v$，便得到一个映射 $T^*:V\to V$。

再证 $T^*$ 线性。对任意标量 $a,b$ 与向量 $v,w$，对一切 $u\in V$ 有
\[
\inner{u}{T^*(av+bw)}=\inner{Tu}{av+bw}=\overline{a}\inner{Tu}{v}+\overline{b}\inner{Tu}{w}
=\inner{u}{aT^*v+bT^*w}.
\]
由内积的非退化性得
\[
T^*(av+bw)=aT^*v+bT^*w.
\]
故 $T^*$ 线性。

最后证唯一性。若 $S$ 也是满足
\[
\inner{Tu}{v}=\inner{u}{Sv}
\]
的线性算子，则对一切 $u,v$ 有
\[
\inner{u}{(T^*-S)v}=0.
\]
由内积的非退化性得 $(T^*-S)v=0$ 对一切 $v$ 成立，所以 $T^*=S$。
\end{proof}

\begin{definition}
若线性算子 $T$ 满足 $T=T^*$，则称 $T$ 为\emph{自伴算子}。在实内积空间中，这等价于“对称算子”；在复内积空间中，这等价于“Hermite 算子”。
\end{definition}

\begin{proposition}
若 $T$ 是自伴算子，则对任意 $v\in V$，$\inner{Tv}{v}\in\R$。
\end{proposition}

\begin{proof}
由自伴性与共轭对称性，
\[
\inner{Tv}{v}=\inner{v}{Tv}=\overline{\inner{Tv}{v}}.
\]
故该数等于其共轭，必为实数。
\end{proof}

\subsection{自伴算子的谱性质}

\begin{theorem}
设 $T$ 是复内积空间上的自伴算子。若 $Tv=\lambda v$，其中 $v\ne 0$，则 $\lambda\in\R$。
\end{theorem}

\begin{proof}
由 $Tv=\lambda v$，
\[
\inner{Tv}{v}=\lambda\inner{v}{v}.
\]
另一方面，由自伴性，
\[
\inner{Tv}{v}=\inner{v}{Tv}=\overline{\lambda}\inner{v}{v}.
\]
因为 $\inner{v}{v}>0$，故 $\lambda=\overline{\lambda}$，即 $\lambda\in\R$。
\end{proof}

\begin{theorem}
设 $T$ 是自伴算子。若 $Tv=\lambda v$，$Tw=\mu w$，且 $\lambda\ne\mu$，则 $v\perp w$。
\end{theorem}

\begin{proof}
由自伴性，
\[
\lambda\inner{v}{w}=\inner{Tv}{w}=\inner{v}{Tw}=\mu\inner{v}{w}.
\]
因此
\[
(\lambda-\mu)\inner{v}{w}=0.
\]
由 $\lambda\ne\mu$ 得 $\inner{v}{w}=0$。
\end{proof}

\begin{lemma}
设 $T$ 是自伴算子，$W\subseteq V$ 是 $T$-不变子空间，即 $T(W)\subseteq W$。则 $W^{\perp}$ 也是 $T$-不变的。
\end{lemma}

\begin{proof}
任取 $x\in W^{\perp}$ 和 $w\in W$。由于 $T(W)\subseteq W$，有 $Tw\in W$。因此
\[
\inner{Tx}{w}=\inner{x}{Tw}=0.
\]
这说明 $Tx$ 与 $W$ 中每个向量都正交，故 $Tx\in W^{\perp}$。
\end{proof}

\begin{theorem}[自伴算子的谱定理：复情形]
设 $V$ 是有限维复内积空间，$T:V\to V$ 是自伴算子。则 $V$ 存在一组由 $T$ 的特征向量组成的正交归一基。换言之，$T$ 在某个正交归一基下的矩阵是实对角矩阵。
\end{theorem}

\begin{proof}
对 $n=\dim V$ 作归纳。若 $n=1$，结论显然。

设 $n>1$。由于复数域上特征多项式必分裂，$T$ 至少有一个特征值 $\lambda\in\C$，并存在非零向量 $v$ 使 $Tv=\lambda v$。由上一定理，$\lambda\in\R$。把 $v$ 归一化后设 $\|v\|=1$，令
\[
W=v^{\perp}.
\]
由上引理，$W$ 对 $T$ 不变。又 $\dim W=n-1$，故按归纳假设，$W$ 存在一组由 $T|_W$ 的特征向量组成的正交归一基 $e_2,\dots,e_n$。令 $e_1=v$，则
\[
e_1,e_2,\dots,e_n
\]
构成 $V$ 的一组正交归一基，而且每个都是 $T$ 的特征向量。

因此在该基下，$T$ 的矩阵是对角矩阵；其对角元正是这些特征值，而它们全是实数。
\end{proof}

\begin{corollary}[实对称矩阵的正交对角化]
设 $A\in M_n(\R)$ 是实对称矩阵，则存在正交矩阵 $Q$ 使得
\[
Q^TAQ=\diag(\lambda_1,\dots,\lambda_n),
\]
其中 $\lambda_i\in\R$。
\end{corollary}

\begin{proof}
把 $A$ 视为复内积空间 $\C^n$ 上的自伴算子。由谱定理可得一组复正交归一特征向量基。由于对应实特征值的特征空间由实线性方程组 $(A-\lambda I)x=0$ 给出，所以每个特征空间都可取实基；再在各特征空间内做 Gram--Schmidt 过程，可得到 $\R^n$ 中的正交归一特征向量基。以这些向量为列组成的矩阵 $Q$ 即所求。
\end{proof}

\begin{theorem}[谱分解]
设 $T$ 是有限维内积空间上的自伴算子，不同特征值为 $\lambda_1,\dots,\lambda_r$，对应特征空间为 $E_{\lambda_i}$。则
\[
V=E_{\lambda_1}\oplus\cdots\oplus E_{\lambda_r},
\]
且这个直和是正交直和。若 $P_i$ 是到 $E_{\lambda_i}$ 上的正交投影，则
\[
T=\sum_{i=1}^r \lambda_i P_i.
\]
\end{theorem}

\begin{proof}
由谱定理，$V$ 存在由 $T$ 的特征向量组成的正交归一基。按特征值分组即可得到
\[
V=E_{\lambda_1}\oplus\cdots\oplus E_{\lambda_r}.
\]
不同特征值的特征空间彼此正交，这是前面定理的直接推论。

任取 $v\in V$，写成
\[
v=v_1+\cdots+v_r\qquad (v_i\in E_{\lambda_i}).
\]
则 $P_i(v)=v_i$，而
\[
Tv=\lambda_1v_1+\cdots+\lambda_rv_r=\sum_{i=1}^r \lambda_i P_i(v).
\]
由于该式对任意 $v$ 成立，故
\[
T=\sum_{i=1}^r \lambda_i P_i.
\]
\end{proof}

\begin{theorem}[交换自伴算子的同时对角化]
设 $T_1,\dots,T_m$ 是有限维复内积空间 $V$ 上两两交换的自伴算子，即
\[
T_iT_j=T_jT_i\qquad (1\le i,j\le m).
\]
则存在 $V$ 的一组正交归一基，使得每个 $T_i$ 在这组基下都表示为实对角矩阵。等价地，$V$ 存在一组公共正交归一本征向量基。
\end{theorem}

\begin{proof}
对 $\dim V$ 归纳。若 $\dim V=1$，结论显然。

先对 $T_1$ 应用谱定理，得
\[
V=\bigoplus_{\lambda} E_{\lambda}(T_1),
\]
其中直和对不同特征值 $\lambda$ 取遍。因为各算子交换，若 $v\in E_{\lambda}(T_1)$，则
\[
T_1(T_jv)=T_j(T_1v)=T_j(\lambda v)=\lambda T_jv,
\]
所以 $T_jv\in E_{\lambda}(T_1)$。这表明每个 $E_{\lambda}(T_1)$ 都对所有 $T_j$ 不变。

把每个 $T_j$ 限制到 $E_{\lambda}(T_1)$ 上，它们仍然两两交换且自伴。由于这些子空间维数严格小于 $\dim V$，根据归纳假设，在每个 $E_{\lambda}(T_1)$ 中都可以找到一组公共正交归一本征向量基。把这些基合并起来，就得到 $V$ 的一组公共正交归一本征向量基。
\end{proof}

\subsection{正算子与正定算子}

\begin{definition}
设 $T$ 是内积空间上的自伴算子。
\begin{enumerate}[label=\textnormal{(\arabic*)}]
  \item 若对任意 $v\in V$，都有 $\inner{Tv}{v}\ge 0$，则称 $T$ 为\emph{正算子}；
  \item 若对任意非零 $v\in V$，都有 $\inner{Tv}{v}>0$，则称 $T$ 为\emph{正定算子}。
\end{enumerate}
\end{definition}

\begin{proposition}
设 $T$ 是自伴算子，则以下命题等价：
\begin{enumerate}[label=\textnormal{(\arabic*)}]
  \item $T$ 是正算子；
  \item $T$ 的全部特征值都非负。
\end{enumerate}
同理，以下命题等价：
\begin{enumerate}[label=\textnormal{(\arabic*)},start=3]
  \item $T$ 是正定算子；
  \item $T$ 的全部特征值都严格为正。
\end{enumerate}
\end{proposition}

\begin{proof}
由谱定理，可取 $V$ 的一组正交归一特征向量基 $e_1,\dots,e_n$，满足
\[
Te_i=\lambda_i e_i,
\]
其中 $\lambda_i\in\R$。

若 $T$ 是正算子，则
\[
\lambda_i=\inner{Te_i}{e_i}\ge 0,
\]
故全部特征值非负。若 $T$ 是正定算子，则同理得 $\lambda_i>0$。

反过来，若全部特征值非负，任取 $v=\sum_i c_ie_i$，则
\[
\inner{Tv}{v}=\sum_i \lambda_i|c_i|^2\ge 0,
\]
故 $T$ 为正算子。若全部特征值严格为正，则对非零 $v$ 至少有一个 $c_i\ne0$，故
\[
\inner{Tv}{v}=\sum_i \lambda_i|c_i|^2>0,
\]
所以 $T$ 正定。
\end{proof}

\begin{corollary}[正平方根]
若 $T$ 是正算子，则存在唯一的正算子 $S$ 使得 $S^2=T$。若 $T$ 正定，则 $S$ 也正定。
\end{corollary}

\begin{proof}
取谱分解
\[
T=\sum_{i=1}^r \lambda_iP_i\qquad (\lambda_i\ge 0).
\]
定义
\[
S=\sum_{i=1}^r \sqrt{\lambda_i}\,P_i.
\]
则 $S$ 自伴，且
\[
S^2=\sum_{i=1}^r \lambda_iP_i=T.
\]
显然 $S$ 是正算子；当所有 $\lambda_i>0$ 时，$S$ 也是正定的。

若 $S'$ 也是正算子且 $(S')^2=T$，则 $S'$ 与 $T$ 交换，并保持 $T$ 的每个特征空间不变。在 $E_{\lambda_i}$ 上，$(S')^2=\lambda_i\Id$，而正性迫使 $S'|_{E_{\lambda_i}}=\sqrt{\lambda_i}\,\Id$。故 $S'=S$，唯一性成立。
\end{proof}

\section{Hermite 矩阵与酉矩阵}

\subsection{矩阵语言中的伴随与正规性}

\begin{definition}
设 $A=(a_{ij})\in M_n(\C)$。定义它的\emph{共轭转置}为
\[
A^*=\overline{A}^{\,t}=(\overline{a_{ji}}).
\]
进一步：
\begin{enumerate}[label=\textnormal{(\arabic*)}]
  \item 若 $A^*=A$，则称 $A$ 为\emph{Hermite 矩阵}；
  \item 若 $A^*A=I$（等价地 $AA^*=I$），则称 $A$ 为\emph{酉矩阵}；
  \item 若 $A\in M_n(\R)$ 且 $A^TA=I$，则称 $A$ 为\emph{正交矩阵}；
  \item 若 $A^*A=AA^*$，则称 $A$ 为\emph{正规矩阵}。
\end{enumerate}
\end{definition}

\begin{proposition}
在标准 Hermite 内积下，矩阵 $A$ 所表示的线性算子的伴随，恰由矩阵 $A^*$ 表示。
\end{proposition}

\begin{proof}
设 $x,y\in\C^n$ 为列向量，则
\[
\inner{Ax}{y}=(Ax)^*y=x^*A^*y=\inner{x}{A^*y}.
\]
因此伴随算子的矩阵就是 $A^*$。
\end{proof}

\begin{proposition}
对 $U\in M_n(\C)$，以下条件等价：
\begin{enumerate}[label=\textnormal{(\arabic*)}]
  \item $U$ 是酉矩阵；
  \item 对任意 $x,y\in\C^n$，有 $\inner{Ux}{Uy}=\inner{x}{y}$；
  \item $U$ 的列向量构成 $\C^n$ 的一组正交归一基；
  \item $U^{-1}=U^*$。
\end{enumerate}
\end{proposition}

\begin{proof}
(1)$\Rightarrow$(2)：若 $U^*U=I$，则
\[
\inner{Ux}{Uy}=x^*U^*Uy=x^*y=\inner{x}{y}.
\]

(2)$\Rightarrow$(3)：令 $u_1,\dots,u_n$ 为 $U$ 的列向量，则
\[
\inner{u_i}{u_j}=\inner{Ue_i}{Ue_j}=\inner{e_i}{e_j},
\]
故列向量正交归一。

(3)$\Rightarrow$(1)：若列向量正交归一，则 $U^*U$ 的 $(i,j)$ 元为 $\inner{u_i}{u_j}$，因此 $U^*U=I$。

(1)$\Leftrightarrow$(4)：由 $U^*U=I$ 立得 $U^*=U^{-1}$；反之亦然。
\end{proof}

\begin{corollary}
若 $U$ 是酉矩阵，则对任意 $x\in\C^n$，有 $\|Ux\|=\|x\|$。若 $Q$ 是正交矩阵，则它保持 $\R^n$ 中的长度与夹角。
\end{corollary}

\begin{proof}
由上命题的 (2) 取 $y=x$ 立即可得。
\end{proof}

\subsection{正规算子的谱定理}

\begin{lemma}
设 $T$ 是有限维复内积空间上的正规算子，即 $T^*T=TT^*$。则对任意 $v\in V$ 与任意 $\lambda\in\C$，有
\[
\|(T-\lambda\Id)v\|=\|(T^*-\overline{\lambda}\Id)v\|.
\]
\end{lemma}

\begin{proof}
直接计算两边平方：
\begin{align*}
\|(T-\lambda\Id)v\|^2
&=\inner{(T-\lambda\Id)v}{(T-\lambda\Id)v}\\
&=\inner{(T^*-\overline{\lambda}\Id)(T-\lambda\Id)v}{v}\\
&=\inner{(T^*T-\lambda T^*-\overline{\lambda}T+|\lambda|^2\Id)v}{v},
\end{align*}
而
\begin{align*}
\|(T^*-\overline{\lambda}\Id)v\|^2
&=\inner{(T^*-\overline{\lambda}\Id)v}{(T^*-\overline{\lambda}\Id)v}\\
&=\inner{(T-\lambda\Id)(T^*-\overline{\lambda}\Id)v}{v}\\
&=\inner{(TT^*-\overline{\lambda}T-\lambda T^*+|\lambda|^2\Id)v}{v}.
\end{align*}
由于 $T^*T=TT^*$，两式相等，因此结论成立。
\end{proof}

\begin{corollary}
若 $T$ 正规，则对任意 $\lambda\in\C$，
\[
\Ker(T-\lambda\Id)=\Ker(T^*-\overline{\lambda}\Id).
\]
特别地，若 $Tv=\lambda v$，则 $T^*v=\overline{\lambda}v$。
\end{corollary}

\begin{proof}
若 $(T-\lambda\Id)v=0$，则由上引理可得
\[
0=\|(T-\lambda\Id)v\|=\|(T^*-\overline{\lambda}\Id)v\|,
\]
故 $(T^*-\overline{\lambda}\Id)v=0$。反向包含同理可得，于是两核相等。
\end{proof}

\begin{theorem}[正规算子的谱定理]
设 $V$ 是有限维复内积空间，$T:V\to V$ 是正规算子。则 $V$ 存在一组由 $T$ 的特征向量组成的正交归一基。等价地，$T$ 在某个正交归一基下的矩阵是对角矩阵。
\end{theorem}

\begin{proof}
仍对 $n=\dim V$ 作归纳。若 $n=1$，结论显然。

由于特征多项式在 $\C$ 上分裂，$T$ 至少有一个特征值 $\lambda$，取单位特征向量 $v$ 满足 $Tv=\lambda v$。由上推论，
\[
T^*v=\overline{\lambda}v.
\]
令 $W=v^{\perp}$。对任意 $w\in W$，有
\[
\inner{Tw}{v}=\inner{w}{T^*v}=\inner{w}{\overline{\lambda}v}=\lambda\inner{w}{v}=0.
\]
故 $Tw\in W$，即 $W$ 对 $T$ 不变。同理也可知 $W$ 对 $T^*$ 不变，因此 $T|_W$ 仍是正规算子。

由归纳假设，$W$ 存在一组正交归一特征向量基 $e_2,\dots,e_n$。令 $e_1=v$，则 $e_1,\dots,e_n$ 构成 $V$ 的一组正交归一特征向量基。
\end{proof}

\begin{corollary}
设 $A\in M_n(\C)$。则以下命题等价：
\begin{enumerate}[label=\textnormal{(\arabic*)}]
  \item $A$ 是正规矩阵；
  \item 存在酉矩阵 $U$ 与对角矩阵 $D$，使得
  \[
  A=UDU^*.
  \]
\end{enumerate}
\end{corollary}

\begin{proof}
(1)$\Rightarrow$(2) 由正规算子的谱定理立即得到。

(2)$\Rightarrow$(1) 直接计算：
\[
A^*A=(UDU^*)^*(UDU^*)=UD^*DU^*,
\]
而
\[
AA^*=UDD^*U^*.
\]
因为对角矩阵彼此交换，所以 $D^*D=DD^*$，故 $A^*A=AA^*$。
\end{proof}

\begin{corollary}
\begin{enumerate}[label=\textnormal{(\arabic*)}]
  \item Hermite 矩阵可被酉矩阵对角化，且对角元全为实数；
  \item 酉矩阵可被酉矩阵对角化，且对角元的模都等于 $1$；
  \item 正交矩阵看作复矩阵后也是酉矩阵，因此可在 $\C^n$ 中酉对角化。
\end{enumerate}
\end{corollary}

\begin{proof}
(1) Hermite 矩阵是自伴矩阵，应用自伴谱定理即可，且特征值实。

(2) 酉矩阵是正规矩阵。若 $Uv=\lambda v$ 且 $v\ne0$，则
\[
\|v\|=\|Uv\|=\|\lambda v\|=|\lambda|\,\|v\|,
\]
故 $|\lambda|=1$。于是由正规谱定理知其可酉对角化，且对角元模长为 $1$。

(3) 正交矩阵满足 $Q^TQ=I$，视为复矩阵即 $Q^*Q=I$，因此是酉矩阵。
\end{proof}

\begin{example}
矩阵
\[
F=\frac{1}{\sqrt2}\mat{1&1\\1&-1}
\]
满足 $F^*F=I$，所以是酉矩阵（也是正交矩阵）。它把标准基变为一组新的正交归一基。另一方面，平面旋转矩阵
\[
R_\theta=\mat{\cos\theta&-\sin\theta\\ \sin\theta&\cos\theta}
\]
在 $\R^2$ 中一般不可对角化；但把它看作复矩阵时，它是酉矩阵，特征值为 $e^{i\theta}$ 与 $e^{-i\theta}$，因此在 $\C^2$ 中可酉对角化。这说明“正交对角化”与“酉对角化”之间确实有域的差别。
\end{example}

\section{与模形式的联系：Petersson 内积}

本节说明本章线性代数如何进入模形式理论。为简洁起见，我们只讨论全模群
\[
\Gamma=\SL_2(\Z)
\]
以及偶权尖点形式空间 $\Sk_k=\Sk_k(\Gamma)$。

\subsection{Petersson 内积的定义}

对 $\gamma=\mat{a&b\\c&d}\in\Gamma$，定义权 $k$ 的 slash 作用
\[
(f|_k\gamma)(z)=(cz+d)^{-k}f\!\left(\frac{az+b}{cz+d}\right).
\]
若 $f\in \Sk_k$，则对每个 $\gamma\in\Gamma$ 都有 $f|_k\gamma=f$。

记标准基本域为
\[
\mathcal{F}=\left\{z=x+iy\in\HH:|x|\le \frac12,\ |z|\ge 1\right\}.
\]

\begin{definition}
设 $f,g\in \Sk_k$。定义它们的\emph{Petersson 内积}为
\[
\inner{f}{g}_{\mathrm{Pet}}=\int_{\mathcal{F}} f(z)\overline{g(z)}\,y^k\,\frac{dx\,dy}{y^2}.
\]
\end{definition}

\begin{lemma}
设
\[
d\mu(z)=\frac{dx\,dy}{y^2}
\]
为上半平面上的双曲测度。则对任意 $\gamma=\mat{a&b\\c&d}\in\SL_2(\R)$，有
\[
d\mu(\gamma z)=d\mu(z),\qquad \Im(\gamma z)=\frac{y}{|cz+d|^2}.
\]
因此若 $f,g\in \Sk_k$，则函数
\[
z\longmapsto f(z)\overline{g(z)}y^k d\mu(z)
\]
对 $\Gamma$ 作用不变。
\end{lemma}

\begin{proof}
记 $z=x+iy$。Mobius 变换的虚部公式是经典恒等式：
\[
\Im\!\left(\frac{az+b}{cz+d}\right)=\frac{\det\gamma\cdot y}{|cz+d|^2}.
\]
当 $\det\gamma=1$ 时即得所述公式。双曲测度的不变性可由直接求 Jacobian 得到，也可从上式推得：Mobius 变换是双曲几何的等距变换，因此保持面积元 $dx\,dy/y^2$。

再看积分核。因为 $f,g$ 都满足模变换律，
\[
f(\gamma z)=(cz+d)^k f(z),\qquad g(\gamma z)=(cz+d)^k g(z).
\]
于是
\begin{align*}
f(\gamma z)\overline{g(\gamma z)}\Im(\gamma z)^k d\mu(\gamma z)
&=(cz+d)^k\overline{(cz+d)^k}f(z)\overline{g(z)}\frac{y^k}{|cz+d|^{2k}}d\mu(z)\\
&=f(z)\overline{g(z)}y^k d\mu(z).
\end{align*}
故积分核对 $\Gamma$ 不变。
\end{proof}

\begin{theorem}
Petersson 内积在 $\Sk_k$ 上良定义，并给出一个 Hermite 内积。
\end{theorem}

\begin{proof}
先证积分收敛。尖点形式在无穷远点有 Fourier 展开
\[
f(z)=\sum_{n\ge 1} a_n e^{2\pi i nz},\qquad g(z)=\sum_{n\ge 1} b_n e^{2\pi i nz}.
\]
在带状区域 $|x|\le 1/2$、$y\ge 1$ 上，级数绝对一致收敛，因此存在常数 $C_f,C_g>0$ 使得
\[
|f(z)|\le C_f e^{-2\pi y},\qquad |g(z)|\le C_g e^{-2\pi y}.
\]
故当 $y\ge 1$ 时，
\[
|f(z)\overline{g(z)}|y^k\frac{1}{y^2}
\le C_fC_g e^{-4\pi y}y^{k-2}.
\]
右端在 $[1,\infty)$ 上可积，因此 cusp 处的积分收敛；而基本域的其余部分是紧的，积分核连续，所以整体积分绝对收敛。

接着验证内积公理。对第一变量线性、对第二变量共轭线性由积分的线性性质立即得到；共轭对称性来自
\[
\overline{\inner{f}{g}_{\mathrm{Pet}}}=\int_{\mathcal{F}} g(z)\overline{f(z)}y^k d\mu(z)=\inner{g}{f}_{\mathrm{Pet}}.
\]

最后证正定性。显然
\[
\inner{f}{f}_{\mathrm{Pet}}=\int_{\mathcal{F}} |f(z)|^2 y^k d\mu(z)\ge 0.
\]
若该积分为零，则连续非负函数 $|f(z)|^2 y^k$ 在 $\mathcal{F}$ 上积分为零。若 $f$ 不恒为零，则存在 $z_0\in\mathcal{F}^{\circ}$ 使 $f(z_0)\ne 0$。连续性给出 $z_0$ 的一个邻域，在该邻域中 $|f(z)|^2 y^k$ 有正下界，于是积分应严格为正，矛盾。故 $f\equiv 0$。因此 Petersson 配对是正定的。
\end{proof}

\begin{remark}
更一般地，对任意有限指数子群 $\Gamma'\subseteq \SL_2(\Z)$ 与尖点形式空间 $S_k(\Gamma')$，同样可以定义 Petersson 内积。其核心原因仍然是：双曲测度与权 $k$ 因子恰好匹配，从而使积分核对群作用不变。
\end{remark}

\subsection{Hecke 算子的自伴性}

设 $\alpha\in \GL_2^+(\Q)$。定义
\[
\alpha^{\vee}=\det(\alpha)\alpha^{-1}.
\]
若双陪集分解为
\[
\Gamma\alpha\Gamma=\bigsqcup_{j=1}^r \Gamma\alpha_j,
\]
则定义作用在 $\Sk_k$ 上的双陪集算子
\[
[\Gamma\alpha\Gamma]f=\det(\alpha)^{k/2-1}\sum_{j=1}^r f|_k\alpha_j.
\]
这是良定义的，因为更换左陪集代表元只会乘以某个 $\gamma\in\Gamma$，而 $f|_k\gamma=f$。

特别地，对
\[
\alpha_n=\mat{1&0\\0&n},
\]
相应的双陪集算子记为 $T_n$，这就是权 $k$ 的 Hecke 算子。它与常见的显式公式
\[
(T_nf)(z)=n^{k-1}\sum_{ad=n\atop d>0}\ \sum_{b\,\mathrm{mod}\,d} d^{-k} f\!\left(\frac{az+b}{d}\right)
\]
是一致的。

\begin{lemma}
对任意 $f,g\in \Sk_k$ 与任意 $\alpha\in \GL_2^+(\Q)$，有
\[
\inner{f|_k\alpha}{g}_{\mathrm{Pet}}=\inner{f}{g|_k\alpha^{\vee}}_{\mathrm{Pet}}.
\]
\end{lemma}

\begin{proof}
这是 Petersson 内积与 slash 作用的基本换元公式。把左边写成在基本域上的积分，并作变量代换 $w=\alpha z$。由于双曲测度对 $\GL_2^+(\R)$ 作用不变，且
\[
\Im(\alpha z)=\frac{\det(\alpha)\,y}{|cz+d|^2},
\]
再结合权 $k$ 的因子，可得变换后恰好出现 $g|_k\alpha^{\vee}$。直接计算便得到
\[
\int_{\mathcal{F}} (f|_k\alpha)(z)\overline{g(z)}y^k d\mu(z)
=
\int_{\mathcal{F}} f(z)\overline{(g|_k\alpha^{\vee})(z)}y^k d\mu(z),
\]
即所求等式。
\end{proof}

\begin{theorem}
对任意 $\alpha\in \GL_2^+(\Q)$，双陪集算子的伴随满足
\[
[\Gamma\alpha\Gamma]^*=[\Gamma\alpha^{\vee}\Gamma].
\]
特别地，对全模群的 Hecke 算子 $T_n$，有
\[
T_n^*=T_n.
\]
\end{theorem}

\begin{proof}
设
\[
\Gamma\alpha\Gamma=\bigsqcup_{j=1}^r \Gamma\alpha_j,
\qquad
\Gamma\alpha^{\vee}\Gamma=\bigsqcup_{j=1}^r \Gamma\beta_j.
\]
根据上一引理，对任意 $f,g\in \Sk_k$，
\begin{align*}
\inner{[\Gamma\alpha\Gamma]f}{g}_{\mathrm{Pet}}
&=\det(\alpha)^{k/2-1}\sum_{j=1}^r \inner{f|_k\alpha_j}{g}_{\mathrm{Pet}}\\
&=\det(\alpha)^{k/2-1}\sum_{j=1}^r \inner{f}{g|_k\alpha_j^{\vee}}_{\mathrm{Pet}}.
\end{align*}
而集合 $\{\alpha_j^{\vee}\}$ 正好给出双陪集 $\Gamma\alpha^{\vee}\Gamma$ 的一组左陪集代表元，因此上式等于
\[
\inner{f}{[\Gamma\alpha^{\vee}\Gamma]g}_{\mathrm{Pet}}.
\]
故
\[
[\Gamma\alpha\Gamma]^*=[\Gamma\alpha^{\vee}\Gamma].
\]

现在取 $\alpha=\alpha_n=\mat{1&0\\0&n}$。则
\[
\alpha_n^{\vee}=\mat{n&0\\0&1}.
\]
设
\[
S=\mat{0&-1\\1&0}\in \Gamma.
\]
直接计算得
\[
S\alpha_nS^{-1}=\alpha_n^{\vee}.
\]
因此
\[
\Gamma\alpha_n\Gamma=\Gamma\alpha_n^{\vee}\Gamma.
\]
由上式便得 $T_n^*=T_n$。
\end{proof}

\subsection{本征模形式分解}

\begin{theorem}
空间 $\Sk_k$ 在 Petersson 内积下是有限维复内积空间。Hecke 算子族 $\{T_n\}_{n\ge 1}$ 两两交换，且每个 $T_n$ 都是自伴算子。因此存在 $\Sk_k$ 的一组正交归一基，由同时满足
\[
T_nf=\lambda_n(f)f\qquad (n\ge 1)
\]
的公共本征向量组成。
\end{theorem}

\begin{proof}
有限维性是模形式基本理论的结论。上一小节已经证明每个 $T_n$ 关于 Petersson 内积自伴；Hecke 算子的交换关系则来自 Hecke 代数的双陪集乘法法则。于是它们构成一族两两交换的自伴算子。

因此可以应用本章已经证明的“交换自伴算子的同时对角化定理”，得到 $\Sk_k$ 的一组公共正交归一本征向量基。每个基向量都是所有 Hecke 算子的公共本征向量，即本征模形式。
\end{proof}

\begin{remark}
在模形式理论中，人们通常再加上 Fourier 展开的首项归一化条件 $a_1(f)=1$，从而得到\emph{归一本征形式}。谱定理给出的正交分解是许多数论结论的起点：Fourier 系数的乘法关系、$L$-函数的 Euler 乘积、Ramanujan 型估计等，都建立在 Hecke 算子的本征分解之上。
\end{remark}

\section*{习题}

下面共 $60$ 题，分为基础题、进阶题与挑战题三组。标记 \basic、\medium、\hard{} 分别表示难度等级。

\subsection*{基础题}

\begin{enumerate}[label=\textbf{5.\arabic*.}]
  \item \basic 在 $\R^2$ 上定义
  \[
  \inner{(x_1,x_2)}{(y_1,y_2)}=2x_1y_1+3x_2y_2.
  \]
  验证这确实是一个内积，并求向量 $(1,-1)$ 的长度。

  \item \basic 在 $\R^2$ 上定义
  \[
  \inner{(x_1,x_2)}{(y_1,y_2)}=x_1y_1-x_2y_2.
  \]
  说明它不是内积，并指出哪一条公理失败。

  \item \basic 在 $\C^2$ 的标准 Hermite 内积下，计算
  \[
  \inner{(1+i,2)}{(i,1-i)}
  \]
  以及向量 $(1+i,2)$ 的范数。

  \item \basic 设
  \[
  A=\mat{a&0\\0&4}.
  \]
  对哪些实数 $a$，公式 $\inner{x}{y}=x^tAy$ 给出 $\R^2$ 上的内积？

  \item \basic 在 $\R^3$ 中计算向量 $u=(1,2,2)$ 与 $v=(2,0,1)$ 的夹角。

  \item \basic 在多项式空间 $P_2(\R)$ 上取内积
  \[
  \inner{f}{g}=\int_0^1 f(x)g(x)\,dx.
  \]
  计算 $\inner{1+x}{1-x}$。

  \item \basic 判断下列向量组是否正交：
  \[
  (1,1,0),\ (1,-1,0),\ (0,0,1).
  \]
  若是，再判断它是否正交归一。

  \item \basic 对向量 $(1,1)$、$(1,-1)$ 应用 Gram--Schmidt 过程，构造 $\R^2$ 的一组正交归一基。

  \item \basic 对向量
  \[
  (1,0,1),\ (1,1,0),\ (0,1,1)
  \]
  应用 Gram--Schmidt 过程。

  \item \basic 将向量组
  \[
  (2,0,0),\ (0,-3,0),\ (0,0,4)
  \]
  归一化为正交归一组。

  \item \basic 证明：在内积空间中，两两正交且非零的向量组一定线性无关。

  \item \basic 求子空间
  \[
  U=\\mathrm{span}\left\{\mat{1\\1\\0},\mat{1\\0\\1}\right\}\subseteq \R^3
  \]
  的一组正交归一基。

  \item \basic 求矩阵
  \[
  A=\mat{1&1\\1&-1}
  \]
  的 QR 分解。

  \item \basic 求矩阵
  \[
  A=\mat{1&0\\1&1\\0&1}
  \]
  的 QR 分解。

  \item \basic 在 $\R^3$ 中，求直线
  \[
  U=\\mathrm{span}\{(1,2,2)\}
  \]
  的正交补 $U^{\perp}$。

  \item \basic 在 $\R^4$ 中，求子空间
  \[
  U=\\mathrm{span}\{(1,0,1,0),(0,1,1,1)\}
  \]
  的正交补的一组基。

  \item \basic 求向量 $v=(2,1,3)$ 在直线 $\\mathrm{span}\{(1,1,0)\}$ 上的正交投影。

  \item \basic 设 $U\subseteq \R^3$ 由正交归一向量
  \[
  e_1=\frac{1}{\sqrt2}(1,1,0),\qquad e_2=(0,0,1)
  \]
  张成。求 $v=(1,2,3)$ 在 $U$ 上的正交投影。

  \item \basic 在 $P_2(\R)$ 上取积分内积，求多项式 $x^2$ 在子空间 $\\mathrm{span}\{1,x\}$ 上的正交投影。

  \item \basic 在 $\R^2$ 中求点 $(3,1)$ 到直线 $y=2x$ 的最佳逼近点。

  \item \basic 用最小二乘法求解超定方程组
  \[
  \begin{cases}
  x=1,\\
  x+y=2,\\
  y=2.
  \end{cases}
  \]

  \item \basic 设
  \[
  A=\mat{2&1\\1&2}.
  \]
  证明它在标准内积下对应自伴算子，并求其特征值与一组正交归一特征向量。

  \item \basic 判断矩阵
  \[
  U=\frac{1}{\sqrt2}\mat{1&1\\-1&1}
  \]
  是否酉；若矩阵元素都视为实数，再判断它是否正交。

  \item \basic 判断矩阵
  \[
  A=\mat{1&i\\-i&2}
  \]
  是否 Hermite。

  \item \basic 将矩阵
  \[
  A=\mat{2&1\\1&2}
  \]
  正交对角化。
\end{enumerate}

\subsection*{进阶题}

\begin{enumerate}[label=\textbf{5.\arabic*.},start=26]
  \item \medium 证明 Cauchy--Schwarz 不等式，并写出等号成立的充要条件。

  \item \medium 从 Cauchy--Schwarz 不等式推出三角不等式。

  \item \medium 证明：若内积空间中的非零向量 $u,v$ 满足
  \[
  |\inner{u}{v}|=\|u\|\,\|v\|,
  \]
  则 $u,v$ 线性相关。

  \item \medium 证明 Gram--Schmidt 过程中构造出来的向量组与原向量组生成相同的子空间。

  \item \medium 证明 QR 分解在要求 $R$ 的对角元为正时是唯一的。

  \item \medium 设 $U$ 是有限维内积空间 $V$ 的子空间。证明
  \[
  (U^{\perp})^{\perp}=U.
  \]

  \item \medium 设 $P_U$ 是到子空间 $U$ 上的正交投影。证明 $P_U$ 线性、满足 $P_U^2=P_U$，并且 $P_U$ 是自伴算子。

  \item \medium 从最佳逼近定理出发，重新推导最小二乘法的正规方程
  \[
  A^*Ax=A^*b.
  \]

  \item \medium 证明：对任意矩阵 $A\in M_{m\times n}(\C)$，矩阵 $A^*A$ 总是正算子；并证明它正定当且仅当 $A$ 的列向量线性无关。

  \item \medium 证明：自伴算子的所有特征值都为实数。

  \item \medium 证明：自伴算子对应不同特征值的特征向量必正交。

  \item \medium 设 $T$ 为自伴算子，$W$ 为 $T$-不变子空间。证明 $W^{\perp}$ 也是 $T$-不变的。

  \item \medium 设 $T$ 是自伴算子，特征值为 $\lambda_1,\dots,\lambda_r$，对应正交投影为 $P_1,\dots,P_r$。证明
  \[
  T=\lambda_1P_1+\cdots+\lambda_rP_r.
  \]

  \item \medium 证明：自伴算子正定当且仅当它的全部特征值都严格为正。

  \item \medium 设 $T$ 是正定算子。证明存在唯一的正定算子 $S$ 使得 $S^2=T$。

  \item \medium 证明：矩阵 $U\in M_n(\C)$ 是酉矩阵，当且仅当它的列向量构成 $\C^n$ 的一组正交归一基。

  \item \medium 证明：若 $U$ 是酉矩阵，则 $|\det U|=1$；若 $Q$ 是实正交矩阵，则 $\det Q=\pm1$。

  \item \medium 证明：酉矩阵的全部特征值都落在单位圆上。

  \item \medium 证明：矩阵 $A$ 正规当且仅当对任意 $x$ 都有
  \[
  \|Ax\|=\|A^*x\|.
  \]

  \item \medium 设 $A$ 是正规矩阵。证明对任意 $\lambda\in\C$，
  \[
  \Ker(A-\lambda I)=\Ker(A^*-\overline{\lambda}I).
  \]

  \item \medium 证明：若一个上三角复矩阵是正规矩阵，则它必为对角矩阵。

  \item \medium 证明：一族两两交换的自伴算子可以同时正交对角化。

  \item \medium 设 $Q\in M_2(\R)$ 为正交矩阵。证明 $Q$ 必形如
  \[
  \mat{\cos\theta&-\sin\theta\\ \sin\theta&\cos\theta}
  \quad\text{或}\quad
  \mat{\cos\theta&\sin\theta\\ \sin\theta&-\cos\theta}
  \]
  其中 $\theta\in\R$。

  \item \medium 设
  \[
  A=\mat{a&b\\ \overline{b}&d}
  \]
  为 $2\times 2$ Hermite 矩阵。证明 $A$ 正定当且仅当 $a>0$ 且 $ad-|b|^2>0$。

  \item \medium 用最小二乘法求经过三点 $(0,1)$、$(1,2)$、$(2,2)$ 的最佳一次拟合多项式，并用正交投影的语言解释结果。
\end{enumerate}

\subsection*{挑战题}

\begin{enumerate}[label=\textbf{5.\arabic*.},start=51]
  \item \hard 设 $f,g\in \Sk_k$。利用它们在无穷远点的 Fourier 展开，严格证明 Petersson 内积
  \[
  \inner{f}{g}_{\mathrm{Pet}}=\int_{\mathcal{F}} f(z)\overline{g(z)}y^k\frac{dx\,dy}{y^2}
  \]
  绝对收敛。

  \item \hard 证明：Petersson 内积与基本域的选取无关。也就是说，若 $\mathcal{F}'$ 是 $\Gamma$ 的另一基本域，则
  \[
  \int_{\mathcal{F}} f(z)\overline{g(z)}y^k d\mu(z)
  =
  \int_{\mathcal{F}'} f(z)\overline{g(z)}y^k d\mu(z).
  \]

  \item \hard 证明：对任意 $\gamma\in\Gamma$，函数
  \[
  f(z)\overline{g(z)}y^k d\mu(z)
  \]
  在变换 $z\mapsto \gamma z$ 下保持不变。

  \item \hard 设 $\alpha\in \GL_2^+(\Q)$。证明换元公式
  \[
  \inner{f|_k\alpha}{g}_{\mathrm{Pet}}=\inner{f}{g|_k\alpha^{\vee}}_{\mathrm{Pet}}.
  \]

  \item \hard 利用上一题，证明双陪集算子的伴随满足
  \[
  [\Gamma\alpha\Gamma]^*=[\Gamma\alpha^{\vee}\Gamma].
  \]

  \item \hard 对 Hecke 算子 $T_n$，证明 $\Gamma\alpha_n\Gamma=\Gamma\alpha_n^{\vee}\Gamma$，从而推出 $T_n$ 关于 Petersson 内积是自伴的。

  \item \hard 证明：由于 Hecke 算子彼此交换且自伴，空间 $\Sk_k$ 存在一组公共正交归一本征模形式基。

  \item \hard 证明：每个有限维内积空间都是完备的。提示：先在一组正交归一基下把它与 $\R^n$ 或 $\C^n$ 等同起来。

  \item \hard 在 Hilbert 空间 $\ell^2(\N)$ 上定义右移算子
  \[
  S(x_1,x_2,x_3,\dots)=(0,x_1,x_2,\dots).
  \]
  证明 $S$ 是等距算子但不是酉算子，也不是自伴算子。

  \item \hard 在 Hilbert 空间 $L^2[0,1]$ 上定义乘法算子
  \[
  (Mf)(x)=xf(x).
  \]
  证明 $M$ 是自伴算子，但没有非零特征向量。由此说明：无限维情形中的谱定理与有限维矩阵对角化并不完全相同。
\end{enumerate}

